\section{Chaleurs spécifiques massique et molaire (diapos 73-76, dont 75-77)}

\begin{objectif}
\textbf{Objectif :} relier la chaleur spécifique $C_v$ et $C_p$ (que l'on utilise tout le temps dans les
exercices) à l'énergie interne $U$ et à l'enthalpie $H$.
\end{objectif}

\textbf{Définitions} : $c=\dfrac{1}{m}\cdot\dfrac{dQ}{dT}$ [J/(kg·K)] (chaleur spécifique \textbf{massique})
et $C=\dfrac{1}{n}\cdot\dfrac{dQ}{dT}$ [J/(mol·K)] (chaleur spécifique \textbf{molaire}), avec $C=M\cdot c$
($M$ = masse molaire). \emph{En mots :} c'est la quantité de chaleur qu'il faut apporter pour élever d'1 K
la température d'1 kg (ou d'1 mole) de matière.

\begin{demotable}
1\textsuperscript{er} principe : $dU=dQ+dW=dQ-p\,dV \Rightarrow dQ=dU+p\,dV$ &
On part du 1\textsuperscript{er} principe pour exprimer $dQ$ (on utilise $dW=-p\,dV$, réversible). \\
$C=\dfrac{1}{n}\cdot\dfrac{dQ}{dT}=\dfrac{1}{n}\left(\dfrac{dU}{dT}+p\dfrac{dV}{dT}\right)$ &
On remplace $dQ$ par l'expression trouvée. C'est la formule \textbf{générale}, valable pour n'importe
quelle transformation. \\
Volume constant : $\boxed{C_v=\dfrac{1}{n}\cdot\dfrac{dU}{dT}}$ &
Si $V$ est constant, $dV=0$ : le deuxième terme disparaît complètement. \\
Pression constante : $C_p=\dfrac{1}{n}\left(\dfrac{dU}{dT}+p\dfrac{dV}{dT}+V\dfrac{dp}{dT}\right)$ &
\textbf{Astuce du cours} : on \emph{ajoute} le terme $V\dfrac{dp}{dT}$, qui vaut $0$ puisque $p$ est
constant ($dp=0$) -- cela ne change donc rien à l'équation, mais permet de compléter une dérivée de produit. \\
$=\dfrac{1}{n}\left(\dfrac{dU}{dT}+\dfrac{d(pV)}{dT}\right)=\dfrac{1}{n}\cdot\dfrac{d(U+pV)}{dT}$ &
On reconnaît la règle du produit $\dfrac{d(pV)}{dT}=p\dfrac{dV}{dT}+V\dfrac{dp}{dT}$, ce qui permet de
regrouper les deux termes. \\
$\boxed{C_p=\dfrac{1}{n}\cdot\dfrac{dH}{dT}}$ &
$U+pV=H$ par définition de l'enthalpie. \\
\end{demotable}

\textbf{Dans l'autre sens} (utile pour les calculs) : de $C_v=\frac{1}{n}\frac{dU}{dT}$ on tire directement
$dU=nC_v\,dT$ -- \emph{cette relation est toujours vraie pour un gaz parfait}, quelle que soit la
transformation, car $U$ ne dépend que de $T$. À volume constant, $W=0$, donc en plus
$\boxed{dQ=dU=nC_v\,dT}$. De même $dH=nC_p\,dT$ est toujours vrai, et à pression constante $W_{ut}=0$
donc $\boxed{dQ=dH=nC_p\,dT}$.

\begin{demotable}
\textbf{Relation de Mayer} : $C_p=\dfrac{1}{n}\dfrac{dU}{dT}+\dfrac{1}{n}\dfrac{d(pV)}{dT}$
$=C_v+\dfrac{1}{n}\dfrac{d(nRT)}{dT}$ &
On repart de la définition de $C_p$ et on utilise $pV=nRT$ (gaz parfait) pour le deuxième terme. \\
$=C_v+\dfrac{1}{n}\cdot n\cdot R$ &
$n$ et $R$ sont des constantes, la dérivée de $nRT$ par rapport à $T$ vaut simplement $nR$. \\
$\boxed{C_p-C_v=R}$ (ou $c_p-c_v=R/M$ en massique) &
On isole. C'est la relation de Mayer, dans le \textbf{formulaire}. \\
$C_p>C_v$, on pose $\dfrac{C_p}{C_v}=\gamma>1$ ($\gamma=1{,}4$ pour l'air) &
Puisque $R>0$, $C_p$ est toujours plus grand que $C_v$ : il faut plus de chaleur pour chauffer un gaz à
pression constante (une partie sert à le dilater) qu'à volume constant. \\
\end{demotable}

\newpage
\section{Transformation adiabatique réversible (diapos 77-84, dont 80-83)}

\begin{objectif}
\textbf{Objectif :} trouver la relation entre $p$, $V$ et $T$ pour une compression/détente
\textbf{rapide} (pas le temps d'échanger de la chaleur) et \textbf{réversible}. C'est la formule la plus
utilisée du cours (moteurs, compresseurs, turbines).
\end{objectif}

\subsection*{Démonstration de la 1\textsuperscript{ère} formule : $T\,V^{\gamma-1}=$ constante}

\begin{demotable}
Adiabatique $\Rightarrow dQ=0$, donc $dW=dU$ &
Par définition, une transformation adiabatique n'échange pas de chaleur ; le 1\textsuperscript{er} principe
$dU=dW+dQ$ se réduit à $dU=dW$. \\
Réversible $\Rightarrow dW=-p\,dV$, et gaz parfait $\Rightarrow dU=nC_v\,dT$ &
On utilise le travail élémentaire réversible (section 1) et la relation $dU=nC_v\,dT$ toujours valable pour
un gaz parfait. \\
$nC_v\,dT+p\,dV=0$ &
On combine les deux lignes précédentes ($dU=dW \Rightarrow nC_v\,dT=-p\,dV$, on réarrange). \\
Gaz parfait : $pV=nRT \Rightarrow p=\dfrac{nRT}{V}$, donc $nC_v\,dT+nRT\dfrac{dV}{V}=0$ &
On remplace $p$ pour n'avoir plus que $T$ et $V$ dans l'équation. \\
$\boxed{C_v\dfrac{dT}{T}+R\dfrac{dV}{V}=0}$ &
On divise toute l'équation par $n$ puis par $T$, pour faire apparaître des rapports «logarithmiques»
($dT/T$ et $dV/V$). \\
On divise par $C_v$ et on utilise $\dfrac{R}{C_v}=\dfrac{C_p-C_v}{C_v}=\gamma-1$ &
On veut faire apparaître $\gamma$ ; comme $C_p-C_v=R$ et $C_p/C_v=\gamma$, on a $R/C_v=\gamma-1$. \\
$\dfrac{dT}{T}+(\gamma-1)\dfrac{dV}{V}=0$ &
Équation réécrite uniquement avec $\gamma$. \\
Intégration : $\ln(T)+(\gamma-1)\ln(V)=$ constante &
$\int \frac{dT}{T}=\ln T$ et $\int\frac{dV}{V}=\ln V$ ; la somme des deux primitives reste une constante. \\
On applique l'exponentielle des deux côtés &
$e^{\ln T + (\gamma-1)\ln V}=e^{\text{constante}}$, ce qui simplifie les logarithmes. \\
$\boxed{T\cdot V^{\gamma-1}=\text{constante}}$ &
Résultat final : c'est la 1\textsuperscript{ère} formule de l'adiabatique réversible. \\
\end{demotable}

\subsection*{Comment obtenir les 2 autres formules (par substitution, avec $pV=nRT$)}

\begin{demotable}
En fonction de $p$ et $V$ : on remplace $T=\dfrac{pV}{nR}$ dans $T\,V^{\gamma-1}=$ cste &
On veut éliminer $T$ au profit de $p$ : on utilise simplement la loi des gaz parfaits pour l'exprimer. \\
$\dfrac{pV}{nR}\cdot V^{\gamma-1}=$ cste $\Rightarrow \dfrac{p\,V^{\gamma}}{nR}=$ cste &
On regroupe les puissances de $V$ ($V\cdot V^{\gamma-1}=V^\gamma$). \\
$n R$ est constant, on l'absorbe dans la constante &
$n$ (quantité de matière) et $R$ ne changent pas pendant la transformation, donc $p\,V^\gamma/(nR)=$cste
$\Leftrightarrow p\,V^\gamma=$ une autre constante. \\
$\boxed{p\,V^\gamma=\text{constante}}$ &
2\textsuperscript{ème} formule, celle utilisée pour lire un diagramme $p$-$V$. \\
\midrule
En fonction de $T$ et $p$ : on remplace $V=\dfrac{nRT}{p}$ dans $T\,V^{\gamma-1}=$ cste &
Même logique, mais on élimine $V$ cette fois. \\
$T\cdot\left(\dfrac{nRT}{p}\right)^{\gamma-1}=$ cste $\Rightarrow T\cdot(nR)^{\gamma-1}\cdot T^{\gamma-1}\cdot
p^{-(\gamma-1)}=$ cste &
On développe la puissance $(\gamma-1)$ sur chaque facteur de la fraction. \\
$T^{\gamma}\cdot p^{-(\gamma-1)}=\dfrac{\text{cste}}{(nR)^{\gamma-1}}$, et $(nR)^{\gamma-1}$ est constant &
On regroupe les puissances de $T$ ($T\cdot T^{\gamma-1}=T^\gamma$) et on absorbe le facteur constant
restant. \\
$\boxed{T^{\gamma}\cdot p^{1-\gamma}=\text{constante}}$ &
3\textsuperscript{ème} formule. Les trois formes sont \textbf{équivalentes} : on choisit celle qui utilise
les variables connues dans l'exercice. \\
\end{demotable}

\begin{retenir}
En dérivant $pV=$cste (isotherme) et $pV^\gamma=$cste (adiabatique) on trouve les pentes
$\frac{dp}{dV}=-\frac{p}{V}$ (isotherme) et $\frac{dp}{dV}=-\frac{p\gamma}{V}$ (adiabatique). Comme
$\gamma>1$, \textbf{une adiabatique est toujours plus pentue qu'une isotherme passant par le même point.}
\end{retenir}
